\section{Результаты}

Программа успешно собрана и запущена в двух операционных системах --- Windows и Linux --- из одной
и той же папки исходных текстов с единой логикой приложений. Реализация обёртки выбиралась
автоматически на этапе конфигурации CMake.

Для проверки вводилась последовательность чисел. Работа подтверждается журналами процессов:
\begin{itemize}
  \item \texttt{[parent pid=]} --- отправленные числа и полученная сумма;
  \item \texttt{[child pid=]} --- полученные числа, нарастающая сумма и факт записи файла-отчёта;
  \item финальный код выхода обоих процессов --- \texttt{0}.
\end{itemize}

\begin{figure}[h]
\centering
\includegraphics[width=0.75\textwidth]{windows.png}
\caption{Запуск \texttt{parent} в Windows}
\end{figure}

\begin{figure}[h]
\centering
\includegraphics[width=0.75\textwidth]{linux.png}
\caption{Запуск \texttt{parent} в Linux}
\end{figure}

Ключевые особенности полученного решения:
\begin{itemize}
  \item кроссплатформенность: код приложений не содержит системных вызовов и полностью идентичен
    для обеих ОС;
  \item корректная передача \texttt{EOF}: закрытие записывающего конца канала родителем завершает
    чтение дочернего процесса;
  \item обработка ошибок: при любой системной ошибке оба процесса завершаются с ненулевым кодом,
    ресурсы освобождаются; на Linux предотвращена гибель родителя по \texttt{SIGPIPE} при записи в
    закрытый канал.
\end{itemize}